\chapter{Tag 19 — Datamatrix komplett von Hand}

\begin{quote}
Greta blickt auf ihren fertigen Encoder. Klassen, Vererbung,
Reed-Solomon, Modi-Wechsler — alles steht. Aber eine Frage lässt sie
nicht los: Wenn jeder Schritt verstanden ist, müsste sie das Symbol
auch \emph{ohne} Computer rechnen können. Drei Codewords, fünf
Prüfbytes, ein 10×10-Raster. Mit Bleistift. Mit Mut.
\end{quote}

\section*{Lernziele für heute}

Am Ende von Tag 19 kannst du \dots

\begin{itemize}
\item GF(256)-Multiplikation mit log/antilog-Tabelle ausführen
\item das Generator-Polynom $g_5(x)$ für Datamatrix konstruieren
\item ein vollständiges 10×10-Datamatrix-Symbol von Hand kodieren —
  Daten, ECC, Platzierung, Frame
\item das eigene Bleistift-Symbol mit dem Smartphone scannen und
  reklamieren: „Da steckt nichts mehr drin, was ich nicht selbst
  gerechnet habe"
\end{itemize}

\paragraph{Material} Bleistift, Radierer, Lineal, kariertes Papier,
Buntstifte (für die UTAH-Markierung). \emph{Kein} Computer für die
Hauptaufgabe. Smartphone für den Scan-Test am Ende. Tabellen und
Vorlagen findest du im Anhang dieses Kapitels — und falls ein Drucker
zur Hand ist, gibt es alle Werkstattbögen einzeln zum Download.

\paragraph{Heute dauert es länger} Die Bleistift-Arbeit braucht
locker drei bis vier Stunden konzentrierte Arbeit, die Bonus-Aufgabe
nochmal so viel. Plane das ein. Es ist der letzte Praktikumstag —
und der ehrlichste.

% =============================================================
\section{Inventur: was steckt noch im Encoder?
         \normalfont(\textit{≈ 10 min})}
% =============================================================

Greta listet auf, was ihr Encoder tut, und prüft jede Komponente
darauf, ob sie sie ohne Computer ausführen könnte:

\begin{center}
\begin{tabular}{lll}
\toprule
\textbf{Schritt} & \textbf{Werkzeug} & \textbf{Hand-Tauglich?} \\
\midrule
Text → ASCII-Codewords         & ASCII-Tabelle    & ja, trivial \\
Daten-Padding                  & PAD-Byte 129     & ja \\
\textbf{GF(256)-Multiplikation} & ?                & \textbf{Hürde} \\
Generator-Polynom $g_5(x)$     & GF(256)-Mul      & wenn Mul geht: ja \\
Reed-Solomon-Polynomdivision   & GF(256)-Mul + XOR & wenn Mul geht: ja \\
Codeword-Platzierung           & Tabelle (Etappe~13) & ja \\
Frame zeichnen                 & Lineal           & ja \\
\bottomrule
\end{tabular}
\end{center}

Die einzige echte Hürde ist die GF(256)-Multiplikation. Die hatten
wir bisher als \texttt{\_gf\_mul()} in Python — sechs Zeilen Bit-Arithmetik
modulo $0x12D$. Per Hand wäre das pro Multiplikation ein 8-Bit-Polynom
mal 8-Bit-Polynom, danach Reduktion. Bei \emph{einer} Multiplikation
machbar, bei \emph{Hunderten} nicht.

\paragraph{Wie viele Multiplikationen?}
Für „MUT" in 10×10: $\sim$15 Multiplikationen für $g_5(x)$ + $3 \times 6 = 18$
Multiplikationen für die RS-Division = ca.\ \textbf{33 GF-Multiplikationen}.
Mit der Logarithmen-Tafel à 30 Sekunden pro Stück sind das knapp
20~Minuten reine Rechenarbeit. Ohne Tafel: pro Multiplikation
mehrere Minuten Bit-Schubsen — nicht vertretbar.

% =============================================================
\section{GF(256) per Logarithmen-Tafel
         \normalfont(\textit{≈ 20 min})}
% =============================================================

Greta erinnert sich an den Mathe-Unterricht der 9.~Klasse: Sinus,
Kosinus und Logarithmus standen damals als Tafelwerk auf dem Pult.
Niemand hat $\sin(37{,}5^\circ)$ aus der Reihenentwicklung gerechnet —
man hat nachgeschlagen. Die Tafel war ein Werkzeug, kein Schummelblatt.

\textbf{In GF(256) funktioniert das genauso.} Statt einer
Sinus-Tafel gibt es eine Logarithmus-Tafel:

\[
  \log[v] = i \quad\text{so dass}\quad \alpha^i = v
\]

mit $\alpha = 2$ und Modulo-Polynom $0x12D$ — die Standard-Wahl
für Datamatrix ECC-200. Die Tafel hat $255$ Einträge (für
$v \in \{1, 2, \dots, 255\}$, der Wert~0 hat keinen Logarithmus).

Dazu eine \textbf{antilog-Tafel}, die rückwärts geht:

\[
  \text{antilog}[i] = \alpha^i \pmod{0x12D}
\]

Mit beiden Tafeln wird Multiplikation zu drei Lookups + einer
Modulo-255-Addition:

\[
  a \times b \;=\; \alpha^{\log[a]} \times \alpha^{\log[b]}
                \;=\; \alpha^{\log[a] + \log[b]}
                \;=\; \text{antilog}\!\big[(\log[a] + \log[b]) \bmod 255\big]
\]

\textbf{Beispiel.} $73 \times 31$ in GF(256):

\begin{enumerate}
\item Nachschlagen: $\log[73] = 198$, $\log[31] = 234$.
\item Addieren: $198 + 234 = 432$.
\item Modulo 255: $432 \bmod 255 = 432 - 255 = 177$.
\item Nachschlagen: $\text{antilog}[177] = 254$.
\item Ergebnis: $73 \times 31 = 254$ in GF(256).
\end{enumerate}

Vier Lookups, eine Addition, eine Modulo-Operation. Das ist der ganze Trick.

\paragraph{Sonderfälle}
\begin{itemize}
\item $a \times 0 = 0$ und $0 \times b = 0$ — kein Tafel-Eintrag nötig.
\item $a \times 1 = a$ — weil $\log[1] = 0$.
\item Die mod-255-Reduktion ist nur nötig, wenn die Summe ≥ 255 wird.
\end{itemize}

% =============================================================
\section{Python-Einheit — die Tafel selbst generieren
         \normalfont(\textit{≈ 15 min})}
% =============================================================

Greta lässt die Tafel einmalig durch Python erzeugen — danach geht es
ohne Computer weiter.

\paragraph{Datei} \texttt{log\_tabelle.py}

\begin{minted}{python}
PRIM = 0x12D     # primitives Polynom für GF(256), Datamatrix ECC-200
ALPHA = 2        # Generator-Element

def baue_tabellen():
    antilog = [0] * 512
    log     = [0] * 256
    x = 1
    for i in range(255):
        antilog[i] = x
        log[x]     = i
        x <<= 1
        if x & 0x100:
            x ^= PRIM
    for i in range(255, 510):
        antilog[i] = antilog[i - 255]   # Bequemlichkeits-Verdoppelung
    return log, antilog

def gf_mul(a, b, log, antilog):
    if a == 0 or b == 0:
        return 0
    return antilog[log[a] + log[b]]
\end{minted}

Die Funktion \texttt{baue\_tabellen()} erzeugt einmalig die zwei
Tabellen. Die antilog-Tabelle ist auf 510~Einträge verdoppelt — so
kann \texttt{gf\_mul} ohne explizite Modulo-Operation arbeiten
(nach dem Lookup landen wir nie über Index~509).

Wer das Skript einmal laufen lässt, kann die Ausgabe abtippen oder
das Buch aufschlagen — die fertigen Tabellen stehen im
\textbf{Anhang zu Tag 19}, gleich auf den nächsten Seiten.%
\footnote{Wer einen Drucker hat, kann den Werkstattbogen
\texttt{zeichenvorlage-gf256.pdf} einzeln herunterladen unter
\url{https://github.com/cgommel/geheimnis-der-redundanz/releases/latest}.}

\begin{bleistiftuebung}{1}{Mit der Tafel multiplizieren}
Schlage die folgenden Multiplikationen mit der log/antilog-Tafel nach:

\begin{enumerate}[label=\alph*)]
\item $78 \times 62$
\item $86 \times 9$
\item $142 \times 30$
\item $228 \times 78$
\item Probe: $5 \times 1 = ?$ Welcher Sonderfall greift?
\end{enumerate}

Die Lösung im Anhang zeigt nicht nur das Ergebnis, sondern jeden Lookup
einzeln. Wenn du dich vertippst, kannst du an jedem Schritt einsteigen.
\end{bleistiftuebung}

% =============================================================
\section{Generator-Polynom $g_5(x)$ für Datamatrix
         \normalfont(\textit{≈ 30 min})}
% =============================================================

Erinnerung aus Etappe 7: das Generator-Polynom für $n_{\text{ecc}}$
Prüfbytes ist
\[
  g_n(x) \;=\; \prod_{i=1}^{n} (x - \alpha^i)
\]
In GF(2) ist Subtraktion gleich Addition gleich XOR — also
$(x - \alpha^i) = (x + \alpha^i)$. Wir bauen $g$ schrittweise auf:

\textbf{Schritt 1:} $g_1(x) = (x + \alpha) = x + 2$.
Koeffizienten (aufsteigend nach Grad): $[2, 1]$.

\textbf{Schritt 2:} $g_2(x) = g_1(x) \cdot (x + \alpha^2)
                            = (x + 2)(x + 4)$.

Multiplikation per Distributivgesetz:
\[
  (x + 2)(x + 4) = x^2 + 4x + 2x + 8 = x^2 + (4 \oplus 2) x + 8
                 = x^2 + 6x + 8
\]

In GF(2)-Addition ist $4 \oplus 2 = 6$ (binär: $100 \oplus 010 = 110$),
nicht~$6$ als Dezimal-Addition — \emph{aber} hier sind beide Wege
gleich. Das ist Zufall: bei höheren Koeffizienten weichen sie ab.
Wer das Symbol $\oplus$ schreibt, vergisst nicht, was eigentlich
gemeint ist.

Koeffizienten: $[8, 6, 1]$.

\textbf{Schritt 3:} $g_3(x) = g_2(x) \cdot (x + \alpha^3) = (x^2 + 6x + 8)(x + 8)$.

\[
  (x^2 + 6x + 8)(x + 8)
  = x^3 + 8x^2 + 6x^2 + 48x + 8x + 64
\]
\[
  = x^3 + (8 \oplus 6) x^2 + (48 \oplus 8) x + 64
  = x^3 + 14 x^2 + 56 x + 64
\]

Hier ist $48 \oplus 8 = 56$ noch dezimal-gleich, aber $8 \oplus 6$:
binär $1000 \oplus 0110 = 1110 = 14$. \emph{Nicht}~$8+6=14$ — die
zwei Wege liefern hier zufällig dasselbe; das ist eine Falle. Greta
notiert \textbf{immer} das XOR-Symbol $\oplus$ und überprüft
binär.

Außerdem: $48 = 6 \cdot 8$ und $64 = 8 \cdot 8$ wurden hier in GF(256)
gerechnet:
\begin{itemize}
\item $6 \times 8$: $\log[6] = 26$, $\log[8] = 3$, $\alpha^{29} = 48$ ✓
\item $8 \times 8$: $\log[8] = 3$, $\alpha^{6} = 64$ ✓
\end{itemize}

Koeffizienten: $[64, 56, 14, 1]$.

\textbf{Schritte 4 und 5} laufen analog mit $\alpha^4 = 16$ bzw.\
$\alpha^5 = 32$. Greta rechnet sie selbst und vergleicht mit dem
Lösungsanhang (Checkpoint~B), der jeden Zwischenstand zeigt.

\textbf{Endergebnis} (siehe Anhang):
\[
  g_5(x) = x^5 + 62 x^4 + 111 x^3 + 15 x^2 + 48 x + 228
\]
Koeffizienten aufsteigend: $[228, 48, 15, 111, 62, 1]$.

% =============================================================
\section{Bleistift-Übung 2: Reed-Solomon für „MUT"
         \normalfont(\textit{≈ 60 min})}
% =============================================================

Jetzt der Hauptakt. Greta nimmt drei Daten-Codewords und rechnet die
fünf ECC-Bytes von Hand.

\paragraph{Schritt 1 — Daten-Codewords}

\[
  \texttt{M} \to 77 + 1 = 78,\quad
  \texttt{U} \to 85 + 1 = 86,\quad
  \texttt{T} \to 84 + 1 = 85
\]

Daten-CW: $[78, 86, 85]$. Eine 10×10-Datamatrix fasst exakt
3~Daten + 5~ECC = 8~Codewords. Kein Padding nötig.

\paragraph{Schritt 2 — RS-Polynomdivision}

Das Schema (aus Etappe 7, jetzt in GF(256)):

\begin{enumerate}
\item Lege ein Rest-Polynom an: $\texttt{rest} = \text{daten} + n_{\text{ecc}}$
  Nullen, also Länge~$3 + 5 = 8$:
  \[
    \texttt{rest} = [78, 86, 85, 0, 0, 0, 0, 0]
  \]
\item Für jedes Daten-Byte $i \in \{0, 1, 2\}$:
  \begin{itemize}
  \item Setze $f = \texttt{rest}[i]$ (der \textbf{Multiplikator}).
  \item Wenn $f \ne 0$: für jedes Glied $j$ von $g_5$ ($j = 0, \dots, 5$):
    \[
      \texttt{rest}[i + j] \;\mathrel{\oplus}\!=\; f \times g_5\text{-Koeffizient}_{[5 - j]}
    \]
    (also: höchster Grad zuerst).
  \end{itemize}
\item Nach drei Iterationen stehen die fünf ECC-Bytes in
  $\texttt{rest}[3 \dots 7]$.
\end{enumerate}

\paragraph{Erste Iteration ($i=0$, $f = 78$)}

Der Multiplikator $f = 78$ wird mit jedem $g_5$-Koeffizienten
multipliziert (höchster Grad zuerst):

\begin{center}
\begin{tabular}{rllr}
\toprule
$j$ & $g$-Koeffizient & GF-Multiplikation & Produkt \\
\midrule
0 & $g[5] = 1$   & $1 \times 78 = $        & $78$ \\
1 & $g[4] = 62$  & $62 \times 78 =$ \textit{(rechnen)} & ? \\
2 & $g[3] = 111$ & $111 \times 78 = $      & ? \\
3 & $g[2] = 15$  & $15 \times 78 = $       & ? \\
4 & $g[1] = 48$  & $48 \times 78 = $       & ? \\
5 & $g[0] = 228$ & $228 \times 78 = $      & ? \\
\bottomrule
\end{tabular}
\end{center}

Das jeweilige Produkt wird per XOR auf
$\texttt{rest}[i + j] = \texttt{rest}[0 + j]$ addiert. Nach Iteration 1:
\[
  \texttt{rest} = [0, ?, ?, ?, ?, ?, 0, 0]
\]
Die führende Null bestätigt: $\texttt{rest}[0]$ wurde durch
$f \times 1 = f$ ausgelöscht. Genau das macht Polynomdivision.

\textbf{Volle Werte} stehen im Lösungsanhang als Checkpoint~C, jede
einzelne GF-Multiplikation ausgeschrieben.

\paragraph{Zweite und dritte Iteration}

Greta wiederholt den Vorgang mit $i = 1$ (neuer Multiplikator
$f = \texttt{rest}[1]$) und $i = 2$. Nach drei Iterationen sind die
ersten drei Plätze von \texttt{rest} ausgenullt, die letzten fünf
enthalten die ECC-Bytes:
\[
  \texttt{rest} = [0, 0, 0, e_0, e_1, e_2, e_3, e_4]
\]

Die finalen Codewords:
\[
  [\,78,\, 86,\, 85,\,\ |\ \,e_0,\, e_1,\, e_2,\, e_3,\, e_4\,]
\]
Spoiler-frei: die ersten drei sind die Daten, die fünf danach kommen
aus der Polynomdivision.

% =============================================================
\section{Platzierung und Frame zeichnen
         \normalfont(\textit{≈ 30 min})}
% =============================================================

Greta hat acht Codewords. Jedes wird in 8 Bits zerlegt (Bit~1 = MSB,
Bit~8 = LSB) und dann nach dem Algorithmus aus Etappe~13 in das
10×10-Symbol einsortiert.

\paragraph{Bit-Zerlegung} Pro Codeword acht Spalten in der
Werkstatt-Tabelle ausfüllen:
\[
  78 = 64 + 8 + 4 + 2 = 0\,1\,0\,0\,1\,1\,1\,0
\]
Hochwertigstes Bit (128er-Stelle) zuerst.

\paragraph{Platzierung} Im Anhang siehst du die innere 8×8-Fläche
mit der vollständigen Slot-Beschriftung
\(\bigl[(\text{CW-Nummer},\,\text{Bit-Nummer})\bigr]\) und einer
Pastellfarbe pro UTAH. Greta malt das Raster auf Karopapier ab
(1~Modul = 1~Karo) und trägt das jeweilige Bit ein — schwarz für~1,
leer für~0.

\paragraph{Frame} Der Frame folgt einem festen Schema:

\begin{itemize}
\item Untere Zeile, linke Spalte: durchgehend schwarz (L-Pattern)
\item Obere Zeile: alternierend, beginnend schwarz
\item Rechte Spalte: alternierend, beginnend weiß
\end{itemize}

Im Anhang ist der Frame bereits eingezeichnet — du musst ihn nur
abmalen.

\begin{bleistiftuebung}{3}{„MUT" auf Karopapier zeichnen}
\begin{enumerate}[label=\alph*)]
\item Übertrage das 10×10-Raster aus dem Anhang auf Karopapier
  (1~Modul = 1~Karo). Frame schon einzeichnen, UTAHs bunt
  einfärben — das hilft beim Zickzack durch das Symbol.
\item Zerlege deine acht Codewords in je 8 Bits. Bei Selbstzweifeln:
  Checkpoint~E im Lösungsanhang zeigt alle 64 Bits.
\item Übertrage die Bits auf dein gezeichnetes Raster anhand der
  Slot-Beschriftung „CW-Nr.\,/\,Bit-Nr.".
\item Scanne dein fertiges Symbol mit dem Smartphone. Es muss „MUT"
  liefern.
\item Falls nicht: vergleiche dein Bitmap mit Checkpoint~G — Modul
  für Modul. Reparaturen mit Bleistift und Radiergummi.
\end{enumerate}
\end{bleistiftuebung}

\paragraph{Wenn der Scan klappt} Du hast einen Datamatrix-Code
\textit{vollständig von Hand} erzeugt. Es gibt nichts mehr im
Encoder, was du nicht selbst gerechnet hast — kein Schritt war Magie.
Wenn dein Smartphone das Symbol liest, gibt dir die kommerzielle
Welt zurück, was du in Bleistift verfasst hast: dein Wort, „MUT",
durch GF(256), Reed-Solomon und ISO/IEC~16022 hindurch, fehlertolerant.

% =============================================================
\section{Bonus: „GESCHAFFT" in 14×14
         \normalfont(\textit{≈ 2--3 h, Hausaufgabe})}
% =============================================================

Wer den 10×10-Code geschafft hat und Lust auf eine Königsdisziplin
hat: derselbe Vorgang in größer.

\paragraph{Aufgabe} „GESCHAFFT" (9 Zeichen) in einem 14×14-Symbol —
mit C40-Modus statt ASCII, weil 9 Zeichen sonst nicht in 8~Daten-CW
passen würden.

\begin{itemize}
\item C40-Triplets: \textbf{GES}, \textbf{CHA}, \textbf{FFT} — drei
  vollständige Triplets, kein Restzeichen.
\item Daten-CW: $[230, b_1, b_2, b_3, b_4, b_5, b_6, 254]$ —
  Latch + 3 Triplets~à 2~CW + Unlatch = 8~CW exakt.
\item ECC: 10 Prüfbytes mit $g_{10}(x)$.
\item RS-Iterationen: 8 statt 3 — etwa $135$ GF-Multiplikationen
  insgesamt. Realistisch: 2--3 Stunden konzentriert.
\end{itemize}

Das 14×14-Raster mit Frame und farbiger UTAH-Markierung steht
ebenfalls im Anhang — direkt nach dem 10×10. Wieder gilt: auf
Karopapier abmalen, UTAHs bunt einfärben, Bits eintragen.\footnote{%
Wer doch einen Drucker hat: \texttt{zeichenvorlage-datamatrix-rs.pdf}
(Seite~2) liegt einzeln im Release-Asset unter
\url{https://github.com/cgommel/geheimnis-der-redundanz/releases/latest}
— bei 135 Multiplikationen ist „eine zweite Vorlage zum nochmal
sauber abschreiben" oft die letzte Rettung.}

Der Lösungsanhang führt durch:

\begin{itemize}
\item Checkpoint~A': die drei C40-Triplet-Rechnungen einzeln
\item Checkpoint~B': $g_{10}(x)$, alle 10 Aufbauschritte tabelliert
\item Checkpoint~C': alle 8 RS-Iterationen mit kompletten Multiplikations-Tabellen
\item Checkpoint~D' bis G': Codewords, Bits, Platzierung, Bitmap
\end{itemize}

Wenn dein 14×14-Symbol „GESCHAFFT" scannt, hast du genau das.

% =============================================================
\section{Reflexion und Würdigung \normalfont(\textit{≈ 20 min})}
% =============================================================

Greta legt Bleistift und Werkstattbogen zur Seite. Auf dem Tisch
liegen zwei Datamatrix-Symbole — eines klein, eines etwas größer.
Beide gezeichnet, gescannt, lesbar.

\subsection*{Was am Anfang stand}

In Etappe~1 war die EAN-13-Prüfziffer eine seltsame Quersumme. Greta
hat sie nachgerechnet, ohne zu wissen warum es ausgerechnet die
gewichtete Summe modulo~10 sein musste. Es funktionierte — das war
genug.

\subsection*{Was sie jetzt versteht}

Die Reise hatte eine innere Logik:

\begin{itemize}
\item \textbf{Etappen 1--2 — Fehler erkennen:} Prüfziffern, Hamming-Distanz.
  Eine einzelne falsche Stelle wird sichtbar, mehr nicht.
\item \textbf{Etappen 3--4 — Fehler korrigieren:} Hamming, SECDED, CRC,
  Interleaving. Ein verkippter Bit lässt sich \emph{lokalisieren} und
  reparieren.
\item \textbf{Etappen 5--6 — Die richtige Mathematik:} Endliche Körper
  $\mathrm{GF}(2^n)$ und Polynome darüber. XOR ist Addition; jede
  Polynom-Operation, die du kennst, läuft hier mit anderen
  Definitionen, aber denselben Strukturen.
\item \textbf{Etappen 7--8 — Reed-Solomon:} Die mächtigste Idee des Buches.
  Polynom-Werte als Codewords, Lagrange-Interpolation als
  Korrekturalgorithmus. Plötzlich kann man $t$~beliebige Fehler in
  $2t$~ECC-Symbolen reparieren — und es gibt eine geschlossene Theorie
  dafür, kein Trick.
\item \textbf{Etappen 9--10 — EAN-13 von Hand und in Pillow:} Strichbreiten,
  Sub-Codierung, Wachzonen. Hardware trifft Theorie.
\item \textbf{Etappen 11--14 — Datamatrix Stück für Stück:} ASCII-Encoding,
  Reed-Solomon-Prüfbytes, ISO-Annex-F-Platzierung, dein erster eigener
  scannbarer Code.
\item \textbf{Etappe 15 — Klassen und Vererbung:} der Code wurde reif. Aus
  einem Skript wurde eine Hierarchie: \texttt{Barcode} oben, mehrere
  Spezialisierungen darunter.
\item \textbf{Etappe 16 — Decoder:} der inverse Pfad. Inverse Platzierung,
  Reed-Solomon-Fehlerlokalisation, ein matplotlib-Stress-Test, der
  zeigt, wo die theoretische Grenze beginnt — und endet.
\item \textbf{Etappe 17 — Große Codes und Interleaving:} ab 36×36 wird das
  Symbol auf mehrere RS-Blöcke aufgeteilt. Ein lokaler Schaden trifft
  jeden Block nur ein bisschen — alle reparieren sich selbst.
\item \textbf{Etappe 18 — Modi:} Digit-Paar, C40, automatischer
  Moduswechsler. Gleicher Inhalt, kleineres Symbol.
\item \textbf{Etappe 19 — Alles von Hand:} GF(256) als Logarithmen-Tafel,
  Generator-Polynom als Produkt, Reed-Solomon als Polynomdivision,
  Platzierung als Tabelle. Der Encoder ist keine Black Box mehr.
\end{itemize}

\subsection*{Was Codierungstheorie eigentlich ist}

Sie ist die mathematische Antwort auf eine sehr alte Frage: Wenn ein
Kommunikationskanal Fehler einführt, wie viel Redundanz brauchen wir,
und wie schlau müssen wir sie verteilen, damit am anderen Ende noch
ankommt, was abgeschickt wurde?

\textbf{Die Antwort hat zwei Lager:} Auf der einen Seite Shannon
(„wie viel Information passt theoretisch durch?"), auf der anderen
Seite die konstruktiven Codes wie Hamming und Reed-Solomon (\,„wie
baut man konkret einen Code, der dieser Grenze nahekommt?"). Das Buch
hat dich auf die zweite Seite geführt — die Bauseite. Und du hast
nicht nur Bibliotheken benutzt, sondern jeden Schritt selbst
nachvollzogen, in mehreren Repräsentationen: Bleistift auf
Karopapier, Python, Klassenhierarchie, Logarithmen-Tafel.

\subsection*{Was bleibt}

Wenn du in zehn Jahren auf einem Pharmakarton ein Datamatrix-Symbol
siehst — auf einem ICE-Sitzplatz, auf einer DHL-Sendung, im
Maschinenbau-Lager, auf einer Pille — wirst du wissen, was darin
steckt. Nicht nur „eine Nummer", sondern: ASCII-Modus oder C40 oder
Digit-Paar; Reed-Solomon mit so-und-soviel Prüfbytes; Multi-Region
oder Single-Region; Interleaving ja oder nein; und, wenn man sehr
genau hinsieht, ein hausgemachtes Bauteil, das du im Praktikum von
Hand gerechnet hast — auch wenn es da gerade als Maschinen-Druck steht.

Ich hoffe, das war den Aufwand wert. Du hast jetzt ein Werkzeug,
das du nie mehr verlierst.

\subsection*{Praktikum vorbei}

Werkstattbogen zusammenrollen. Bleistift in den Becher. Datamatrix-Symbol
mit „MUT" über den Schreibtisch hängen — kleines Trophäen-Stück.
Wenn jemand fragt, wofür das ist, sag: \glqq Das ist ein Datamatrix.
Steht „MUT" drin. Hab ich von Hand gerechnet. Reed-Solomon und alles.\grqq{}
Und wenn sie fragen, was Reed-Solomon ist, hast du jetzt sehr viel zu
erzählen.
